% =====================================================================
% 05_prop2.tex — Proposition 2: realized and expected excess returns (WP-C)
% =====================================================================
\section{Proposition 2: realized and expected excess returns}
\label{sec:prop2}

This section proves Proposition~2 of KL: how a flight-to-safety shock
$\dev{\cy}_1>0$ moves the \emph{realized} and \emph{expected} returns on the
three assets in the simplified economy --- safe dollar bonds, capital, and
Foreign bonds. The entire derivation reuses the impulse-response engine of
Section~\ref{sec:engine}; we do \emph{not} re-derive any reduced form. As in
Proposition~\ref{prop:prices}, we specialize to identical portfolios, no Home
government safe debt ($b^g_{H,s}=0$), and complete home bias
$\hbias\to\zs/(1+\zs)$, so that $\dev{\wsh}_1=0$ and $\Eone\dev{\wsh}_2=0$
(established at \eqref{eq:prop1-wsh0}). The labor responses
$(\dev{\ell}_1,\fr{\dev{\ell}}_1)$ are exactly those of
Proposition~\ref{prop:prices}: zero under flexible wages, and
\eqref{eq:prop1b-impact} under predetermined wages. We treat the two wage
regimes in turn for each asset.

\subsection{Statement}

\begin{proposition}[Realized and expected excess returns]
\label{prop:returns}
Consider the simplified economy of Definition~\ref{def:simplified} with identical
portfolios across countries and no Home government safe debt ($b^g_{H,s}=0$), so
that $\dev{\wsh}_1=0$ and $\Eone\dev{\wsh}_2=0$, and consider a flight-to-safety
shock $\dev{\cy}_1>0$ in period~$1$ with $\widetilde{\widetilde{\dev k}}_1=0$ on
impact. Then, to first order, on impact of the shock:
\begin{enumerate}[label=(\alph*),leftmargin=1.8em,itemsep=3pt]
  \item the realized real return on safe dollar bonds \emph{rises},
  $\rr_1\propto\dev{\cy}_1$ ($>0$), in both wage regimes;
  \item the realized real return on capital is \emph{unaffected} if wages are
  flexible, $\rk_1=0$, but \emph{falls} if wages are set one period in advance,
  $\rk_1\propto-\dev{\cy}_1$ ($<0$);
  \item the realized real return on Foreign bonds (in Home consumption units) is
  \emph{unaffected} if wages are flexible, $\fr{r}_1-\D\dev{\rer}_1=0$, but
  \emph{falls} if wages are set one period in advance,
  $\fr{r}_1-\D\dev{\rer}_1\propto-\dev{\cy}_1$ ($<0$);
  \item the \emph{expected} excess returns on capital and on Foreign bonds over
  dollar bonds are positive and equal to the convenience yield:
  \begin{equation}
  \Eone\bigl[\rk_2-\rr_2\bigr]
  =\Eone\bigl[\fr{r}_2-\D\dev{\rer}_2-\rr_2\bigr]
  =\dev{\cy}_1. \label{eq:prop2-wedge}
  \end{equation}
\end{enumerate}
\end{proposition}

\noindent (We transcribe the proposition to match KL's main-paper statement
(p.~1663). The proportionality constants in (a)--(c) are made explicit in the
proof; here $\D\dev{\rer}_t\equiv\dev{\rer}_t-\dev{\rer}_{t-1}$ is the change in
the real exchange rate, and since the economy is at its deterministic steady
state in period~$0$, $\dev{\rer}_0=0$ so $\D\dev{\rer}_1=\dev{\rer}_1$.)

\subsection{Proof}

\paragraph{Specialization of the engine.} As in the proof of
Proposition~\ref{prop:prices} (their app.~B.3.1), identical portfolios and
$b^g_{H,s}=0$ make the realized wealth-share equation~\eqref{eq:wsh-realized}
vanish, and complete home bias $\hbias\to\zs/(1+\zs)$ gives
$\Eone\dev{\wsh}_2=0$ \eqref{eq:prop1-wsh0}, so the composite denominator
collapses to $\Lambda=1/(1-\cshare)$ \eqref{eq:Lambda-cmplt} and the
relative-labor coefficient is $\hbias/\Lambda=\frac{\zs}{1+\zs}(1-\cshare)$. We
set $\widetilde{\widetilde{\dev k}}_1=0$ on impact throughout. Recall the labor
solutions: flexible wages give $\dev{\ell}_1=\fr{\dev{\ell}}_1=0$
\eqref{eq:flex-labor}, while predetermined wages give \eqref{eq:prop1b-impact}
\[
\dev{\ell}_1 = -\frac{1}{\taylor}\,\frac{1}{\cshare+\frac{1}{\taylor}(1-\cshare)}\,\dev{\cy}_1 < 0,
\qquad \fr{\dev{\ell}}_1 = 0.
\]
For compactness write
\begin{equation}
D \;\equiv\; \cshare+\tfrac{1}{\taylor}(1-\cshare) \;>\;0,
\qquad
\dev{\ell}_1 = -\frac{1}{\taylor D}\,\dev{\cy}_1, \label{eq:prop2-Ddef}
\end{equation}
so that under predetermined wages $\dev{\ell}_1<0$ for $\dev{\cy}_1>0$. We also
reuse, from the proof of Proposition~\ref{prop:prices}, the predetermined-wage
values (eqs.~\eqref{eq:prop1b-Erk}, \eqref{eq:prop1b-q})
\begin{equation}
\Eone\rk_2 = -(1-\cshare)\dev{\ell}_1 \;>\;0,
\qquad
\dev{\rer}_1 = \dev{\tot}_1 = -(1-\cshare)\dev{\ell}_1 \;>\;0, \label{eq:prop2-reuse}
\end{equation}
where the real-exchange-rate value uses $\hbias\frac{1+\zs}{\zs}=1$ at complete
home bias (so $\dev{\rer}_1=\dev{\tot}_1$) and $\Lambda=1/(1-\cshare)$.

\subparagraph{(a) Realized return on safe dollar bonds.} The realized dollar-bond
return is minus realized Home inflation, $\rr_1=-\D\dev P_1$
\eqref{eq:realized-bond}, with $\D\dev P_1$ from the Home Taylor
rule~\eqref{eq:dP1}, so $\rr_1=\frac{1}{\taylor}(\dev{\cy}_1-\Eone\rk_2)$ by
\eqref{eq:r1-solved}.

\emph{Flexible wages.} Here $\dev{\ell}_1=\fr{\dev{\ell}}_1=0$, so by
\eqref{eq:Er2k-solved} $\Eone\rk_2=0$, and
\begin{equation}
\rr_1 = -\D\dev P_1 = \frac{1}{\taylor}\dev{\cy}_1 \;>\;0. \label{eq:prop2a-flex}
\end{equation}
This is KL's ``$\rr_1=\frac{1}{\phi}\dev{\cy}_1>0$.''

\emph{Predetermined wages.} Now $\Eone\rk_2=-(1-\cshare)\dev{\ell}_1>0$ by
\eqref{eq:prop2-reuse}, so by \eqref{eq:r1-solved}
\begin{equation}
\rr_1 = \frac{1}{\taylor}\bigl(\dev{\cy}_1-\Eone\rk_2\bigr)
= \frac{1}{\taylor}\bigl(\dev{\cy}_1+(1-\cshare)\dev{\ell}_1\bigr)
= \frac{1}{\taylor}\dev{\cy}_1\left[1-\frac{1-\cshare}{\taylor D}\right]. \label{eq:prop2a-sticky}
\end{equation}
The bracket is
$1-\frac{1-\cshare}{\taylor D}=\frac{\taylor D-(1-\cshare)}{\taylor D}
=\frac{\taylor\cshare+(1-\cshare)-(1-\cshare)}{\taylor D}
=\frac{\taylor\cshare}{\taylor D}=\frac{\cshare}{D}\in(0,1)$, using
$\taylor D=\taylor\cshare+(1-\cshare)$ from \eqref{eq:prop2-Ddef}. Hence
\begin{equation}
\rr_1 = \frac{1}{\taylor}\frac{\cshare}{D}\dev{\cy}_1
=\frac{\cshare}{\taylor\cshare+(1-\cshare)}\dev{\cy}_1 \;>\;0,
\qquad\text{i.e.}\quad \rr_1\propto\dev{\cy}_1. \label{eq:prop2a-sticky2}
\end{equation}
In both regimes $\rr_1\propto\dev{\cy}_1>0$: the realized real return on dollar
bonds rises, establishing claim~(a). (Consistently, $\rr_1=-\D\dev P_1$ and
$\D\dev P_1<0$ in both regimes by Proposition~\ref{prop:prices}: Home deflates,
so the real return on the predetermined nominal bond rises.)

\subparagraph{(b) Realized return on capital.} We read $\rk_1$ directly from the
solved engine reduced form~\eqref{eq:rk1-solved}. With $\Eone\dev{\wsh}_2=0$ and
$\widetilde{\widetilde{\dev k}}_1=0$ the wealth-share and capital terms drop:
\begin{equation}
\rk_1 = -\frac{\hbias}{\Lambda}\bigl(\fr{\dev{\ell}}_1-\dev{\ell}_1\bigr)
+(1-\cshare)\Bigl(\tfrac{1}{1+\zs}\dev{\ell}_1+\tfrac{\zs}{1+\zs}\fr{\dev{\ell}}_1\Bigr). \label{eq:prop2b-rk}
\end{equation}

\emph{Flexible wages.} With $\dev{\ell}_1=\fr{\dev{\ell}}_1=0$, both terms vanish:
\begin{equation}
\rk_1 = 0. \label{eq:prop2b-flex}
\end{equation}

\emph{Predetermined wages.} Using $\fr{\dev{\ell}}_1=0$,
$\hbias/\Lambda=\frac{\zs}{1+\zs}(1-\cshare)$, and
$\frac{1}{1+\zs}\dev{\ell}_1+\frac{\zs}{1+\zs}\fr{\dev{\ell}}_1
=\frac{1}{1+\zs}\dev{\ell}_1$,
\begin{align}
\rk_1 &= -\frac{\zs}{1+\zs}(1-\cshare)\bigl(0-\dev{\ell}_1\bigr)
+(1-\cshare)\frac{1}{1+\zs}\dev{\ell}_1 \notag\\
&= (1-\cshare)\dev{\ell}_1\left[\frac{\zs}{1+\zs}+\frac{1}{1+\zs}\right]
= (1-\cshare)\dev{\ell}_1 \;<\;0, \label{eq:prop2b-sticky}
\end{align}
since $\frac{\zs}{1+\zs}+\frac{1}{1+\zs}=1$ and $\dev{\ell}_1<0$. Substituting
$\dev{\ell}_1=-\frac{1}{\taylor D}\dev{\cy}_1$,
\begin{equation}
\rk_1 = -(1-\cshare)\frac{1}{\taylor D}\dev{\cy}_1 \;\propto\; -\dev{\cy}_1. \label{eq:prop2b-sticky2}
\end{equation}
This establishes claim~(b): $\rk_1=0$ under flexible wages and
$\rk_1\propto-\dev{\cy}_1<0$ under predetermined wages. (Note
$\rk_1=(1-\cshare)\dev{\ell}_1=-\Eone\rk_2$ by \eqref{eq:prop2-reuse}; the
realized capital return moves \emph{opposite} to the expected capital return on
impact, a fact we use in the discussion.)

\subparagraph{(c) Realized excess return on Foreign bonds.} The object of
interest is the realized return on Foreign bonds expressed in Home consumption
units \emph{relative to} the real-exchange-rate change, $\fr{r}_1-\D\dev{\rer}_1$.
Since $\dev{\rer}_0=0$ (steady state), $\D\dev{\rer}_1=\dev{\rer}_1$. The Foreign
realized bond return is $\fr{r}_1=-\D\fr{\dev P}_1$ \eqref{eq:realized-bond}, with
$\D\fr{\dev P}_1$ from \eqref{eq:dPstar1}:
\begin{equation}
\fr{r}_1 = -\D\fr{\dev P}_1
= -\frac{1}{\taylor}\left(\Eone\rk_2-\frac{1+\zs}{\zs}\frac{\hbias}{\Lambda}\bigl(\fr{\dev{\ell}}_1-\dev{\ell}_1\bigr)\right). \label{eq:prop2c-rstar}
\end{equation}

\emph{Flexible wages.} With $\dev{\ell}_1=\fr{\dev{\ell}}_1=0$ we have
$\Eone\rk_2=0$, so $\fr{r}_1=0$; and $\dev{\rer}_1=\dev{\tot}_1=0$ by
\eqref{eq:s1-solved} (both labor terms and $\Eone\dev{\wsh}_2$ vanish), so
$\D\dev{\rer}_1=0$. Hence
\begin{equation}
\fr{r}_1-\D\dev{\rer}_1 = 0-0 = 0. \label{eq:prop2c-flex}
\end{equation}

\emph{Predetermined wages.} Use $\fr{\dev{\ell}}_1=0$, complete home bias
$\frac{1+\zs}{\zs}\frac{\hbias}{\Lambda}=\frac{\hbias\frac{1+\zs}{\zs}}{\Lambda}
=\frac{1}{\Lambda}=(1-\cshare)$ (since $\hbias\frac{1+\zs}{\zs}=1$ and
$\Lambda=1/(1-\cshare)$), and $\Eone\rk_2=-(1-\cshare)\dev{\ell}_1$. Then
\eqref{eq:prop2c-rstar} becomes
\begin{align}
\fr{r}_1 &= -\frac{1}{\taylor}\Bigl(-(1-\cshare)\dev{\ell}_1-(1-\cshare)\bigl(0-\dev{\ell}_1\bigr)\Bigr)
= -\frac{1}{\taylor}\Bigl(-(1-\cshare)\dev{\ell}_1+(1-\cshare)\dev{\ell}_1\Bigr) \notag\\
&= 0. \label{eq:prop2c-rstar-val}
\end{align}
\gapflag{KL state claim~(c) in the main paper but suppress the algebra. The
cancellation in \eqref{eq:prop2c-rstar-val} is exact: with $\fr{\dev{\ell}}_1=0$
the two pieces inside the parenthesis are $\Eone\rk_2=-(1-\cshare)\dev{\ell}_1$
and $\frac{1+\zs}{\zs}\frac{\hbias}{\Lambda}(\fr{\dev{\ell}}_1-\dev{\ell}_1)
=(1-\cshare)(-\dev{\ell}_1)=-(1-\cshare)\dev{\ell}_1$, which are equal, so
$\fr{r}_1=-\D\fr{\dev P}_1=0$. Economically: at complete home bias and
$\fr{\dev{\ell}}_1=0$, Foreign output and Foreign inflation are unchanged on
impact (the safety shock hits only the Home Taylor rule), so the Foreign nominal
bond's realized real return is zero.} Meanwhile, by \eqref{eq:prop2-reuse},
$\D\dev{\rer}_1=\dev{\rer}_1=-(1-\cshare)\dev{\ell}_1>0$. Therefore
\begin{equation}
\fr{r}_1-\D\dev{\rer}_1 = 0-\bigl(-(1-\cshare)\dev{\ell}_1\bigr)
= (1-\cshare)\dev{\ell}_1 \;<\;0,
\qquad \fr{r}_1-\D\dev{\rer}_1\propto-\dev{\cy}_1. \label{eq:prop2c-sticky}
\end{equation}
This establishes claim~(c): the Home-currency realized return on Foreign bonds is
zero under flexible wages and falls under predetermined wages. The driver is the
real \emph{appreciation} of the Home currency ($\D\dev{\rer}_1>0$): a Foreign
bond loses value in Home consumption units when the dollar appreciates. (Note
that $\fr{r}_1-\D\dev{\rer}_1=(1-\cshare)\dev{\ell}_1=\rk_1$ by
\eqref{eq:prop2b-sticky}: the realized excess returns on capital and on Foreign
bonds coincide on impact under these specializations.)

\subparagraph{(d) Expected excess returns.} The expected wedges are read directly
from the linearized Euler block of the engine, \eqref{eq:euler-r} and
\eqref{eq:euler-rstar} (their app.~eqs.~(84)--(85)). The capital wedge
\eqref{eq:euler-r} is, in difference form,
\begin{equation}
\Eone\rk_2 = \Eone\rr_2 + \dev{\cy}_1
\;\Longrightarrow\;
\Eone\bigl[\rk_2-\rr_2\bigr] = \dev{\cy}_1, \label{eq:prop2d-cap}
\end{equation}
which holds \emph{regardless} of the wage regime (it is the household's
intertemporal condition, not a function of the labor solution). This is the first
equality in \eqref{eq:prop2-wedge}.

For the Foreign-bond wedge, the engine's expected Foreign-return relation
\eqref{eq:euler-rstar} reads
\begin{equation}
\Eone\fr{r}_2 = \Eone\rr_2 + \dev{\cy}_1 + \hbias\,\frac{1+\zs}{\zs}\,\D\Eone\dev{\tot}_2. \label{eq:prop2d-rstar-raw}
\end{equation}
Here $\D\Eone\dev{\tot}_2\equiv\Eone\dev{\tot}_2-\dev{\tot}_1$ is the expected
period-$1$-to-$2$ change in the terms of trade. We must show this collapses to
the clean statement $\Eone[\fr{r}_2-\D\dev{\rer}_2-\rr_2]=\dev{\cy}_1$, i.e.\ that
the terms-of-trade adjustment exactly accounts for the expected
real-exchange-rate change $\Eone\D\dev{\rer}_2$. Two facts close this.

By the period-$2$ real-exchange relation~\eqref{eq:Eq2},
$\Eone\dev{\rer}_2=\hbias\frac{1+\zs}{\zs}\Eone\dev{\tot}_2$, and by the
period-$1$ relation~\eqref{eq:q1-s1},
$\dev{\rer}_1=\hbias\frac{1+\zs}{\zs}\dev{\tot}_1$. Differencing the two,
\begin{equation}
\Eone\D\dev{\rer}_2 \equiv \Eone\dev{\rer}_2-\dev{\rer}_1
= \hbias\,\frac{1+\zs}{\zs}\bigl(\Eone\dev{\tot}_2-\dev{\tot}_1\bigr)
= \hbias\,\frac{1+\zs}{\zs}\,\D\Eone\dev{\tot}_2. \label{eq:prop2d-Dq2}
\end{equation}
Substituting \eqref{eq:prop2d-Dq2} into \eqref{eq:prop2d-rstar-raw} gives,
\emph{identically in all parameters and both wage regimes},
\begin{equation}
\Eone\fr{r}_2 = \Eone\rr_2 + \dev{\cy}_1 + \Eone\D\dev{\rer}_2
\;\Longrightarrow\;
\Eone\bigl[\fr{r}_2-\D\dev{\rer}_2-\rr_2\bigr] = \dev{\cy}_1, \label{eq:prop2d-rstar}
\end{equation}
the second equality in \eqref{eq:prop2-wedge}.
\gapflag{KL write \eqref{eq:euler-rstar}/(85) as
$\Eone\fr r_2=\Eone\rr_2+\dev{\cy}_1+\hbias\frac{1+\zs}{\zs}\D\Eone\dev{\tot}_2$
``where we have used~(76)'' (their \eqref{eq:Eq2}), and state the wedge result
$\Eone[\fr r_2-\D\dev{\rer}_2-\rr_2]=\dev{\cy}_1$ in Proposition~2 without
re-displaying the one-line identification of
$\hbias\frac{1+\zs}{\zs}\D\Eone\dev{\tot}_2$ with $\Eone\D\dev{\rer}_2$. We supply
that identification: it is just the period-$2$-minus-period-$1$ difference of the
two real-exchange-rate relations \eqref{eq:q1-s1} and \eqref{eq:Eq2}, both of
which carry the \emph{same} proportionality constant $\hbias\frac{1+\zs}{\zs}$,
so the difference is
$\hbias\frac{1+\zs}{\zs}\D\Eone\dev{\tot}_2=\Eone\D\dev{\rer}_2$ exactly. The
result therefore holds for general $\hbias$, not only at complete home bias; the
home-bias specialization is \emph{not} needed for part~(d). Under the
Proposition's specialization $\Eone\dev{\wsh}_2=0$ one further has
$\Eone\dev{\tot}_2=0$ \eqref{eq:Es2}, so $\D\Eone\dev{\tot}_2=-\dev{\tot}_1$ and
$\Eone\D\dev{\rer}_2=-\dev{\rer}_1$; but the wedge identity
\eqref{eq:prop2d-rstar} does not rely on this, only on the common constant of
proportionality.}

\medskip
Combining \eqref{eq:prop2d-cap} and \eqref{eq:prop2d-rstar} yields the displayed
equation~\eqref{eq:prop2-wedge}:
\[
\Eone\bigl[\rk_2-\rr_2\bigr]
=\Eone\bigl[\fr{r}_2-\D\dev{\rer}_2-\rr_2\bigr]=\dev{\cy}_1>0,
\]
the expected excess returns on capital and on Foreign bonds over dollar bonds are
both positive and equal to the convenience yield, establishing claim~(d). \qed
(\,Proposition~\ref{prop:returns}\,)

\begin{remark}[Realized vs.\ expected: a sign flip]
Collecting the predetermined-wage results, on impact
\[
\rr_1 = \frac{\cshare}{\taylor\cshare+(1-\cshare)}\dev{\cy}_1 > 0,
\qquad
\rk_1 = \fr{r}_1-\D\dev{\rer}_1 = (1-\cshare)\dev{\ell}_1 < 0,
\]
so the \emph{realized} excess returns on capital and Foreign bonds are
\[
\rk_1-\rr_1 < 0,
\qquad
\fr{r}_1-\D\dev{\rer}_1-\rr_1 < 0,
\]
both strictly negative, while the \emph{expected} excess returns going forward
are $\Eone[\rk_2-\rr_2]=\Eone[\fr r_2-\D\dev{\rer}_2-\rr_2]=\dev{\cy}_1>0$. The
risky assets underperform safe dollar bonds exactly on impact of the flight to
safety, and are expected to outperform thereafter. This realized-vs-expected sign
flip is the engine of the dollar's negative risk premium (Proposition~5).
\end{remark}

\subsection{Discussion}

\paragraph{Realized excess returns are negative on impact.} Proposition~2 says
that on impact of a positive safety shock the realized excess returns on capital
and on Foreign bonds are \emph{negative}: both assets pay badly exactly when the
flight to safety hits. Two distinct forces produce this, with different scope.

\emph{Force 1 (deflation; operative even with flexible wages).} The flight to
safety forces Home deflation, $\D\dev P_1<0$ (Proposition~\ref{prop:prices}),
which raises the realized real return on the predetermined nominal dollar bond,
$\rr_1=-\D\dev P_1>0$. This channel is present \emph{even under flexible wages}:
in that case $\rk_1=0$ and $\fr{r}_1-\D\dev{\rer}_1=0$, yet $\rr_1>0$, so the
realized excess returns $\rk_1-\rr_1$ and $\fr{r}_1-\D\dev{\rer}_1-\rr_1$ are
strictly negative purely because the dollar bond's real return rises. The safe
bond is a hedge against safety shocks through deflation alone.

\emph{Force 2 (recession and real appreciation; requires nominal rigidity).} With
wages set in advance, the safety shock additionally causes a Home-concentrated
recession ($\dev{\ell}_1<0$) and a real dollar appreciation ($\dev{\rer}_1>0$).
The recession raises the expected marginal product of capital, so the price of
installed capital falls today ($\dev{\qk}_1<0$), pulling down the realized
capital return: $\rk_1=(1-\cshare)\dev{\ell}_1<0$. And the real appreciation
directly lowers the Home-currency realized return on Foreign bonds:
$\fr{r}_1-\D\dev{\rer}_1=(1-\cshare)\dev{\ell}_1<0$. So under nominal rigidity
both risky assets fall in \emph{absolute} terms, not merely relative to the
appreciating dollar bond.

\paragraph{Expected excess returns are positive going forward.} For agents to
hold capital and Foreign bonds despite their poor impact returns, these assets
must be \emph{expected} to outperform dollar bonds going forward. This is exactly
the Euler-equation wedge \eqref{eq:prop2-wedge}: the convenience yield
$\dev{\cy}_1$ is the compensation, in expected excess return, that safe dollar
bonds forgo because of the safety/liquidity services they provide. A larger
safety shock raises both the impact underperformance of risky assets and the
forward-looking expected excess return that restores indifference.

\paragraph{Connection to the target facts: the dollar as a hedge.} The
realized-vs-expected sign flip is the model's microfoundation for the empirical
fact that \emph{safe dollar bonds are a hedge}: they deliver high realized real
returns precisely in states (flights to safety) when capital and foreign-currency
bonds deliver low returns. An asset that pays well in bad times commands a
\emph{negative} risk premium and has a \emph{negative beta} with respect to the
global risk factor --- the dollar's ``exorbitant privilege'' as a negative-beta
asset. Proposition~2 supplies the realized covariance structure (dollar bonds up,
risky assets down, in the same shock) that Proposition~\ref{prop:wealth} and
Proposition~5 will turn into wealth redistribution and a signed risk premium. In
the big\_cy reading, a \emph{large absolute} convenience yield $\dev{\cy}_1$
magnifies the wedge in \eqref{eq:prop2-wedge} and hence the magnitude of the
dollar's negative risk premium, without requiring any asymmetry in risk tolerance
across countries --- the convenience yield does the work that KL otherwise load
onto greater US risk-bearing capacity.
